\section{Welfare Analysis} \label{sec:mvpf}

Following \cite{hendren2020unified} and \cite{finkelstein2020welfare}, we conduct welfare analysis to evaluate the overall impact of the road maintenance tax cuts and compute a conservative estimate of the Marginal Value of Public Funds (MVPF). We define MVPF as the ratio of household's willingness to pay (WTP) for the tax cut to the government's net cost of providing the tax cut. Specifically, we compute MVPF as follows:

\begin{equation}
    \text{MVPF} = \frac{\text{WTP}}{\text{Net Cost to Government}} = \frac{\text{Tax Savings - House Value Decline}}{\text{Direct Long Run Cost + Fiscal Externality}}
\end{equation}

Where,
\begin{itemize}
    \item \textbf{Direct Long Run Cost} is the tax revenue loss per household to the government due to the road maintenance tax cut over the life of the tax cut.
    \item \textbf{Fiscal Externality} is the additional loss in property tax revenue to the government due to the decline in tax base following the decline in house prices.
    \item \textbf{Tax Savings} is the overall value of tax savings from the road maintenance tax cut.
    \item \textbf{House Value Decline} is the average of the yearly capitalized loss in house values over the period of analysis.  
\end{itemize}

Using this framework, we compute an MVPF of -\$1.4, indicating that for every dollar of revenue loss the government, the households incur a net loss of \$1.4, after accounting for both direct costs and fiscal externalities. This negative MVPF suggests that the road maintenance tax cut leads to a net welfare loss for society, making it a welfare destroying policy. To put this in perspective, although \cite{hendren2020unified} do not report a sub-category for public infrastructure and focus on government spending programs instead of tax cuts, their MVPF estimates for policies targeting adults range from 0.5 to 2.

